\section{Precision Loss in Optimizer Comparisons}
\label{sec:problem}
Figure~\ref{fig:projection-trap} gives the smallest example.  A feature matrix maps several public optimizer contracts into the same label.  After the mapping, behaviors with different placement, action kind, and evidence layer become indistinguishable, so a study that treats the label as equality can compare non-equivalent public optimizer behavior.  Existing benchmarks and optimizer-comparison methods standardize workloads, cardinalities, or plan-quality measurements; they do not define the equality relation over public pushdown, adaptation, or plan-evidence contracts.  OptSem-C fills that gap before a JOB-, TPC-, or engine-specific measurement is interpreted.

\begin{figure}[t]
\centering
\begin{tikzpicture}[
  band/.style={vpanel, draw=vldbRule, fill=vldbSoft, minimum width=7.14cm,
    minimum height=1.06cm},
  card/.style={vpanel, draw=vldbBlue!68, fill=vldbMist, text width=2.10cm,
    minimum height=.92cm, font=\rmfamily\fontsize{6.75}{7.35}\selectfont,
    text=black!72, align=center, inner xsep=2.3pt, inner ysep=1.6pt,
    text height=1.38ex, text depth=.24ex},
  collide/.style={vpanel, draw=vldbRose!76!black, fill=vldbRose!8,
    text width=2.76cm, minimum height=.74cm,
    font=\rmfamily\bfseries\fontsize{6.95}{7.55}\selectfont,
    text=vldbRose!78!black, align=center, inner xsep=2.2pt, inner ysep=1.5pt,
    text height=1.45ex, text depth=.28ex},
  repair/.style={vpanel, draw=vldbTeal!72!black, fill=vldbMint,
    text width=2.10cm, minimum height=.78cm,
    font=\rmfamily\fontsize{6.75}{7.35}\selectfont, text=black!70,
    align=center, inner xsep=2.5pt, inner ysep=1.6pt,
    text height=1.38ex, text depth=.24ex},
  label/.style={font=\rmfamily\bfseries\fontsize{7.15}{7.85}\selectfont,
    text=vldbInk, align=left, inner sep=.3pt, text height=1.45ex,
    text depth=.28ex},
  lost/.style={vpanel, draw=vldbRose!68!black, fill=white, text width=1.18cm,
    minimum height=.72cm, font=\rmfamily\fontsize{6.60}{7.20}\selectfont,
    text=vldbRose!76!black, align=center, inner xsep=1.2pt, inner ysep=1.2pt,
    text height=1.24ex, text depth=.18ex},
  chip/.style={rounded corners=1.3pt, draw=black!14, fill=white,
    font=\rmfamily\fontsize{6.65}{7.25}\selectfont, text=vldbGray,
    inner xsep=1.9pt, inner ysep=.55pt, text height=1.32ex,
    text depth=.25ex},
  flow/.style={-{Latex[length=1.6mm]}, line width=.70pt,
    draw=black!48, line cap=round, line join=round}
]
\path[use as bounding box] (-.05,-4.30) rectangle (7.32,.24);

\begin{scope}[on background layer]
  \node[band] at (3.60,-.72) {};
  \node[band, fill=vldbWarm] at (3.60,-2.08) {};
  \node[band, fill=vldbMint] at (3.60,-3.44) {};
\end{scope}

\node[label, anchor=west] at (.12,-.19) {A. Reference contracts};
\node[label, anchor=west] at (.12,-1.56) {B. Coarse projection};
\node[label, anchor=west] at (.12,-2.92) {C. Repaired keys};

\node[card] (m1) at (1.58,-.86) {\textbf{$M_1$}\\filter\\local plan};
\node[card] (m2) at (5.62,-.86) {\textbf{$M_2$}\\aggregate\\source stage};
\node[lost] (diff) at (3.60,-.86) {different\\operator\\placement};
\node[chip, draw=vldbRose!35, text=vldbRose!70!black] at (3.60,-.30)
  {$S(M_1)\neq S(M_2)$};

\node[collide] (p) at (3.60,-2.12) {\texttt{pushdown=yes}\\declares equal};
\node[chip, draw=vldbRose!35, text=vldbRose!70!black] at (1.56,-2.58)
  {operator dropped};
\node[chip, draw=vldbRose!35, text=vldbRose!70!black] at (5.64,-2.58)
  {placement dropped};

\node[repair] (r1) at (1.58,-3.62) {filter\\+ local};
\node[repair] (r2) at (5.62,-3.62) {aggregate\\+ source};
\node[chip, draw=vldbTeal!45, text=vldbTeal!58!black] (key) at (3.60,-3.22)
  {operator + placement};

\draw[flow, draw=vldbRose!72!black] (m1.south) -- ++(0,-.28) -| (p.north west);
\draw[flow, draw=vldbRose!72!black] (m2.south) -- ++(0,-.28) -| (p.north east);
\draw[flow, draw=vldbTeal!72!black] (p.south) -- (key.north);
\draw[flow, draw=vldbTeal!72!black] (key.west) -- (r1.east);
\draw[flow, draw=vldbTeal!72!black] (key.east) -- (r2.west);
\end{tikzpicture}
\caption{Projection loss as a finite contract witness. Two source-backed contracts become equal under a coarse \texttt{pushdown=yes} denominator; retaining operator and placement separates the witness without adding private optimizer facts.}
\Description{Two reference contracts with different operator and placement fields collapse into a single pushdown=yes projection. Restoring operator and placement separates the contracts into distinct repaired keys.}
\label{fig:projection-trap}
\end{figure}


\textbf{Representative witness audits.}  Table~\ref{tab:cases} gives six concrete witness rows before the formal model.  Each starts from a familiar public label, names the engine pair and probe, shows the contract distinction that the label hides, and gives the comparison rewrite.  C1 is the running P0009 example; C2--C3 cover keyword join/plan evidence, C4--C5 cover operator-only plan and join rows, and C6 is the narrow yes/no boundary.  The same finite checker covers the 492 primary keyword/operator records and the six yes/no boundary records.

\begin{table}[t]
\centering
\caption{Concrete collision audits linking coarse labels to hidden contract pairs. C1 binds BigQuery federated pushdown~\cite{bigquery_federated_docs} and plan-stage evidence~\cite{bigquery_docs} to Trino pushdown~\cite{trino_pushdown_docs} and dynamic-filtering evidence~\cite{trino_dynamic_filtering_docs}.}
\label{tab:cases}
\footnotesize
\setlength{\tabcolsep}{1.6pt}
\renewcommand{\arraystretch}{1.05}
\begin{tabular}{@{}>{\raggedright\arraybackslash}p{0.24in}>{\raggedright\arraybackslash}p{1.82in}>{\raggedright\arraybackslash}p{0.74in}@{}}
\toprule
Case & Hidden contract contrast & Separating fields \\
\midrule
C1 & pushdown/explain: BigQuery delegation+stage evidence vs. Trino pushdown+dynamic filtering & \textbf{variant, placement, decision time} \\
C2 & plan evidence: PostgreSQL estimates vs. BigQuery/Snowflake stages & \textbf{observability, placement} \\
C3 & adaptive: compile-time planning vs. runtime repartition/AQE/conversion & \textbf{layer, decision time} \\
C4 & CTE: DuckDB/PostgreSQL inline, materialize, reuse, and disable guards & \textbf{variant, guard state} \\
\bottomrule
\end{tabular}
\end{table}


C1 is mechanical.  Source spans admit rules; P0009 triggers them with a remote aggregate; reference signatures differ on placement, layer, decision time, and observability; and the keyword projection maps both to \texttt{pushdown}+\texttt{explain}.  A coarse existence claim may remain, but portability or latency explanations must keep those fields.

\textbf{Feature-name collision.} A cell such as \texttt{pushdown=yes} may denote predicate movement inside an engine, partition pruning, connector-side predicate delegation, aggregation delegation, join delegation, limit pushdown, Top-N delegation, or runtime filtering.  These actions affect different operators and have different guards, placements, decision times, and observable plan effects.  Runtime and adaptivity distinctions are visible in Tukwila's adaptive integration~\cite{ives1999tukwila}, dynamic query plans~\cite{cole1994optimization}, and Spark SQL adaptive execution~\cite{armbrust2015spark}.  Rewrite and execution-layer distinctions appear in outerjoin rewrites~\cite{galindo1997outerjoin}, compiled execution~\cite{neumann2011efficient}, Calcite's adapter-centered optimizer framework~\cite{begoli2018calcite}, Orca's modular optimizer architecture~\cite{soliman2014orca}, and Dremel-style serving trees~\cite{melnik2010dremel}.  A single feature name can therefore erase exactly the fields a comparison needs.

\textbf{Operator disappearance.} A missing native plan node is not a reliable semantics.  It may indicate source delegation, logical-plan observability, rewrite, stage-level aggregation, or runtime-only visibility.  OptSem-C represents source-delegated execution as a positive contract action rather than as absence.  This matters for modern analytical systems whose public surfaces combine Snowflake-style cloud-service stages~\cite{dageville2016snowflake}, Redshift's managed warehouse architecture~\cite{gupta2015redshift}, C-Store column layout~\cite{stonebraker2005cstore}, row/column execution differences~\cite{abadi2008columnstores}, column-oriented design guidance~\cite{abadi2013columnoriented}, materialization policy~\cite{abadi2007materialization}, compression-coupled execution~\cite{abadi2006compression}, embedded DuckDB execution~\cite{raasveldt2019duckdb}, vectorized MonetDB/X100 execution~\cite{boncz2005monetdb}, HyPer's compiled in-memory engine~\cite{kemper2011hyper}, Vectorwise vectorized execution~\cite{zukowski2012vectorwise}, Hive's MapReduce warehouse layer~\cite{thusoo2009hive}, Impala's Hadoop SQL engine~\cite{kornacker2015impala}, distributed SQL services such as F1~\cite{shute2013f1}, and dataflow-style execution such as Spark RDDs~\cite{zaharia2012resilient}.

\textbf{Architecture ranking.} A local engine should not be ranked below a distributed engine merely because broadcast, shuffle, or remote-source delegation is not applicable.  Conversely, a federated or cloud engine should not be credited with local optimizer behavior when the public contract exposes delegation or service-side stages.  OptSem-C separates non-applicability, documented absence, documented optional behavior, and unspecified evidence.

\textbf{Observable contract versus hidden implementation.} The object of study is not what an engine could do internally.  It is what the public contract supports as a comparison claim.  This distinction is analogous to the separation between SQL equivalence and optimizer-contract equivalence.  Cosette proves SQL equivalences by reducing query semantics to logic~\cite{chu2017cosette}; semiring-based decision procedures give another equivalence foundation~\cite{chu2018usemiring}; and QED expands practical SQL equivalence checking~\cite{wang2024qed}.  Chandra and Merlin characterize conjunctive-query containment~\cite{chandra1977optimal}; Aho et al. characterize equivalence among relational expressions~\cite{aho1979equivalences}.  OptSem-C asks a different question: whether the public optimizer behavior being compared exposes the same evidence layers and comparison obligations.

These precision losses imply a minimum payload for comparing public optimizer behavior: guard, canonical action, modality, placement, decision time, observability, version, and provenance.  The rest of the paper formalizes this payload and measures what happens when a comparison vocabulary projects it away.
